%% paper2_revised.tex - revised manuscript after the ICACS-2026 rejection.
%% Venue-neutral article format with author-year citations (IJPHM uses APA
%% style; swap in the IJPHM template once the target is confirmed).
%% Every number in this file comes from results/revised/*.json; see
%% revision/08_change_log.md for what changed and why.
\documentclass[11pt]{article}
\usepackage[T1]{fontenc}
\usepackage[margin=1in]{geometry}
\usepackage{amsmath,amssymb}
\usepackage{graphicx}
\usepackage{booktabs}
\usepackage{xcolor}
\usepackage{url}
\usepackage[round,authoryear]{natbib}
\usepackage[hidelinks]{hyperref}
\graphicspath{{figures/}{./}}

\title{Protocol Choices Decide the Winner: A Re-examination of Lightweight
Unsupervised Anomaly Detectors for Industrial Sensor Data, with Retraining
Policies and Resource Costs}
\author{Sarjan Shah \and Abhi Shah\\[2pt]
\small Department of Information Technology, Lok Jagruti University, Vadodara, Gujarat, India\\
\small \texttt{shahsarjan968@gmail.com}, \texttt{abhi.shah@ljku.edu.in}}
\date{}

\begin{document}
\maketitle

\begin{abstract}
Comparisons of unsupervised anomaly detectors for industrial sensor data often
report a winning method, but the winner can depend on evaluation choices that are
rarely varied. This paper re-examines an earlier six-detector comparison on two
public water-system benchmarks, SKAB and SWaT, after an audit found that its SWaT
copy (a public re-upload) placed the attack period's normal rows in the file named
for normal operation, that thresholds were calibrated with different rules for
different detectors, and that its drift adaptation selected retraining data with
ground-truth labels. Nine lightweight detectors and a small LSTM autoencoder are
evaluated under standard protocols with threshold-free metrics (AUROC, AUPRC),
threshold-dependent metrics under two calibration rules, event recall and
block-bootstrap confidence intervals. On SWaT's standard split, all ten detectors
rank attacks equally well (AUROC $0.83$--$0.86$, overlapping intervals); on SKAB
with pooled training, none is reliably better than chance, while SKAB's per-run
protocol yields run-level AUROC up to $0.79$ and an inverted ranking. Protocol choices
moved MCC by up to $0.8$ for the same detector. With thresholds set only from normal
data, the $z$-score, One-Class SVM and Isolation Forest were the most dependable on
SWaT (MCC $0.64$--$0.72$), while LOF, PCA reconstruction and the autoencoder raised
alarms on about 70\% of normal readings in the attack week. In a streaming
simulation, SWaT's normal operation drifted enough to raise a static Isolation
Forest's false-alarm rate from 12\% to 76\%; retraining restored detection only when
retraining data were known to be normal, drift-triggered retraining was no better
than a schedule with the same budget, and label-free retraining failed. Single-reading
cost in a one-CPU, 1\,GB container ranged from microseconds to over 100\,ms. The
results argue for choosing detectors by threshold stability and retraining needs and
for reporting protocols in enough detail to tell which of these a benchmark number
measures.
\end{abstract}

\noindent\textbf{Keywords:} unsupervised anomaly detection; multivariate time
series; evaluation protocol; concept drift; retraining policy; resource-constrained inference

\section{Introduction}

Pumps, valves and treatment plants are instrumented with sensors whose joint
behaviour changes before or during a fault. Because faults and attacks are rare
and costly to label, condition monitoring often relies on \emph{unsupervised}
(novelty) detectors that are fitted to normal operation and flag departures from
it \citep{chandola2009}. For an engineer choosing such a detector for a
resource-limited monitoring node, three practical questions matter: which
detector separates anomalies from normal operation best, whether that answer
survives the passage of time once the detector is deployed, and what the
detector costs to run.

Published comparisons give less guidance on these questions than their
headline numbers suggest. Large benchmark studies report that no family of
methods is consistently best and that deep models often do not beat
well-chosen classical ones \citep{audibert2022,schmidl2022,sarfraz2024}. A
parallel line of work shows that evaluation choices, such as point adjustment,
metrics that reward trivial detectors, and flawed benchmark labels, can
manufacture apparent progress \citep{kim2022,garg2021,wu2023,liu2024}. These
critiques concern mainly metrics and labels. Less attention has been paid to
the other protocol choices an applied study must make: how the training and
test data are defined, how a detection threshold is calibrated, and, once
deployed, how and on which data a detector is retrained.

This paper started as a six-detector comparison on two industrial benchmarks,
the Skoltech Anomaly Benchmark (SKAB) \citep{skab} and the Secure Water
Treatment testbed (SWaT) \citep{goh2016,mathur2016}. An audit of that study
found that its conclusions depended on such protocol choices. The SWaT copy that
was used (a public re-upload) stores the normal-labelled rows of the attack
period in the file named for normal operation, so the original study trained
on the normal intervals \emph{between} attacks rather than on the separate
normal run. Thresholds had been calibrated with different rules for different
models, and the drift-adaptation experiment selected its retraining data using
ground-truth labels. Rather than only correcting these points, the revised
study measures how much each choice changes the conclusions.

The contributions are empirical and methodological, not algorithmic:
\begin{enumerate}
\item A controlled comparison of nine lightweight unsupervised detectors and a
small LSTM autoencoder on SKAB and SWaT under standard, documented protocols,
reported with threshold-free metrics (AUROC, AUPRC), threshold-dependent
metrics (MCC, event recall, false-alarm rate) and block-bootstrap confidence
intervals (Section~\ref{sec:results}).
\item A protocol-sensitivity analysis showing how the split definition, the
training protocol and the threshold rule change the measured performance and
the ranking of detectors (Section~\ref{sec:protocol}).
\item A retraining-policy study on SWaT's genuine normal run and attack period
that compares never, periodic, drift-triggered and continuous retraining, with
retraining data selected with or without labels (Section~\ref{sec:drift}).
\item Single-reading latency, throughput, model size and peak memory for every
detector inside a one-CPU, 1\,GB ARM64 container (Section~\ref{sec:edge}).
\end{enumerate}

The study answers four research questions:
\begin{description}
\item[RQ1] Under standard, documented protocols with uncertainty estimates, do
lightweight detectors differ in how well they rank anomalies on SKAB and SWaT?
\item[RQ2] How much do the split definition, the training protocol and the
threshold rule change measured performance and the ranking of detectors?
\item[RQ3] Does drift-triggered retraining outperform periodic or continuous
retraining, and what happens when retraining data are selected without labels?
\item[RQ4] What does each detector cost per reading under a one-CPU, 1\,GB
memory cap?
\end{description}

\section{Related Work}
\label{sec:related}

\paragraph{Detectors.} Classical unsupervised detectors differ in what they
treat as normal. Distance- and density-based methods such as the Local Outlier
Factor (LOF) score a point by the relative density of its neighbourhood
\citep{breunig2000}; boundary methods such as the One-Class SVM enclose the
bulk of the training data \citep{scholkopf2001}; covariance methods flag large
Mahalanobis distances, with the Minimum Covariance Determinant (MCD) giving a
robust estimate \citep{rousseeuw1999}; PCA reconstruction scores the residual
outside a low-rank normal subspace \citep{shyu2003}; Isolation Forest isolates
points by random partitioning \citep{liu2008}; and the histogram- and
copula-based HBOS, COPOD and ECOD combine per-feature tail probabilities and are
fast and nearly parameter-free \citep{goldstein2016,li2022ecod,zhao2019}. Deep
sequence models, including LSTM encoder-decoders \citep{malhotra2016},
prediction-error detectors \citep{hundman2018}, adversarial autoencoders
\citep{audibert2020}, graph networks \citep{deng2021} and transformers
\citep{tuli2022}, report strong results on SWaT, usually with protocols and
metrics that differ from paper to paper.

\paragraph{What benchmark studies agree on.} Comparisons that put conventional,
machine-learning and deep detectors side by side find that no family wins
consistently \citep{audibert2022}, that multivariate problems are harder on
average \citep{schmidl2022}, and that simple or linear baselines match many
recent deep models once inflated metrics are removed \citep{sarfraz2024}. Tabular
benchmarks reach similar conclusions \citep{goldstein2016,han2022}. A result of
the form ``detector X wins on dataset A but not on B'' is therefore expected
and is not, on its own, a contribution.

\paragraph{Where evaluations go wrong.} Point adjustment credits a whole event
when any of its points is flagged and can make random scores look competitive
\citep{kim2022}. Point-wise F1 and accuracy can reward detectors that flag
almost everything, which the Matthews correlation coefficient (MCC) does not
\citep{chicco2020,garg2021}. Threshold-dependent metrics make comparisons hinge
on the threshold rule; threshold-free AUC measures and their range-aware
extensions avoid this \citep{paparrizos2022,liu2024}, and range-based
precision and recall account for events rather than points \citep{tatbul2018}.
Benchmark data can itself be flawed \citep{wu2023}. These studies examine metrics
and labels; the protocol choices examined here (split semantics, threshold
calibration rule, and the label assumption in retraining) are complementary.

\paragraph{Drift and retraining.} Concept-drift research distinguishes
detecting a change, understanding it and adapting to it \citep{gama2014,lu2018}.
Two-sample tests on feature distributions are a simple and competitive way to
detect shift \citep{rabanser2019}; adaptive windows such as ADWIN resize the
reference automatically \citep{bifet2007}; streaming anomaly benchmarks such as
NAB score detectors online \citep{ahmad2017}. For an anomaly detector, adapting
to drift carries a specific risk: data used for retraining may contain the very
anomalies the detector should flag. Our search found no study that compares
drift-triggered with periodic retraining at matched cost on SWaT or quantifies
the effect of selecting retraining data with versus without labels; this is a
statement about the search performed (documented in the supplementary
literature log), not a claim of absence.

\paragraph{Resource cost.} MLPerf Tiny measures accuracy, latency and energy of
single inferences on constrained hardware, including an anomaly-detection task
\citep{banbury2021}. Comparable cost reporting for classical multivariate
detectors on industrial data is rare. Here, cost is measured as single-reading
latency, batch throughput, model size, load time and peak memory under explicit
CPU and memory caps. Energy is not measured.

\section{Datasets}
\label{sec:data}

\paragraph{SKAB.} SKAB was recorded on a water-circulation testbed (pump,
pipes, valves) with eight sensors sampled at 1\,Hz: two accelerometers (RMS),
motor current, pressure, water temperature, pump-housing temperature, voltage
and volumetric flow rate \citep{skab}. It contains one 9{,}401-row anomaly-free
run and 34 fault runs (46{,}860 rows in total) in three families: valve~1
(16 runs), valve~2 (4) and other faults (14). Faults are physically induced and
labelled per second. Each file is an independent experiment.

\paragraph{SWaT.} SWaT is a six-stage water-treatment plant with 51 sensors and
actuators sampled at 1\,Hz \citep{goh2016,mathur2016}. The standard release has a
seven-day normal run (22--28 December 2015, 496{,}800 rows) and an
attack period (28 December--2 January, 449{,}919 rows) in which 54{,}621 rows
(12.1\%) are labelled as attacks. The copy used in this study is a public
re-upload (Kaggle, \texttt{vishala28/swat-dataset-secure-water-treatment-system}, v3)
in which (i) \texttt{normal.csv} contains the normal run twice and, in its first
395{,}298 rows, the normal-labelled rows of the attack period, and (ii)
\texttt{attack.csv} contains only the 54{,}621 attack rows. In this copy seven
channels are blank throughout the normal run (MV101, AIT201, MV201, P201, P202,
P204, MV303). The standard split was reconstructed by de-duplicating the normal
run by timestamp and merging the attack-period rows of both files in time order,
which reproduces the published row counts. The 24 continuous sensors that are
present in both periods are used (AIT201, a continuous analyser, is excluded
because it is blank in the normal run). Following common practice, the first
21{,}600 rows of the normal run (start-up stabilisation) are discarded
\citep{li2019madgan}. Results should be re-checked on the official iTrust release.

\section{Method}
\label{sec:method}

\subsection{Features}
For each continuous sensor, the raw value, 5-second and 30-second causal
rolling means, a 5-second rolling standard deviation and the first difference
are computed (40 features on SKAB, 120 on SWaT). Features are computed on each
complete recording \emph{before} any rows are removed, so windows never join
readings across a removed fault and, at test time, use only past readings.
The original study's elapsed-time feature was removed: its range differs between
training (up to 9{,}960\,s on SKAB) and test segments (at most 291\,s), and it
has no meaning for a monitor that runs continuously. Its effect is reported
in Section~\ref{sec:protocol}. Wall-clock features remain excluded because they
leak the identity of the recording session.

\subsection{Detectors}
Nine classical detectors are evaluated, all producing a continuous anomaly score:
a maximum absolute $z$-score; PCA reconstruction error (95\% of variance);
Elliptic Envelope (robust MCD on ten whitened principal components, needed
because the features are collinear); Isolation Forest (100 trees); One-Class SVM
(RBF kernel, $\nu{=}0.1$); LOF ($k{=}20$, novelty mode); and ECOD, COPOD and
HBOS from PyOD \citep{zhao2019}, with library defaults. On SWaT, One-Class SVM and
LOF are fitted on a 20{,}000-row evenly spaced subsample of the training data
because their cost grows quadratically; all other detectors use every training
row. The deep baseline is an LSTM autoencoder (one-layer encoder and decoder,
64 hidden units, $\approx$\,52k--58k parameters) on 30-second windows of raw
scaled sensors, trained for 25 epochs with Adam; a row's score is the
reconstruction error of the window ending at that row \citep{malhotra2016}.
All models are fitted on normal rows only.

\subsection{Threshold rules}
\label{sec:thr}
A score becomes an alarm when it exceeds a threshold $\tau$. Because the
comparison should not depend on one rule, two rules are applied identically to
every detector:
\begin{itemize}
\item \textbf{Label-free:} $\tau$ is the 99th percentile of scores on held-out
\emph{normal} data, so about 1\% of normal readings are expected to raise an
alarm. No anomaly labels are used.
\item \textbf{Validation-MCC:} $\tau$ maximises MCC on a labelled validation
set that contains anomalies, excluding thresholds that flag more than 99\% of
rows. This is semi-supervised calibration and is reported as such.
\end{itemize}
A third rule, maximising validation F1, is used only to show how much the
original study's choice of F1 mattered.

\subsection{Splits}
\paragraph{SKAB, pooled protocol.} As in the original study, each fault run is
split chronologically into its first 80\% and last 20\%. The training portions of
20\% of the runs (whole runs, fixed seed) form a labelled validation set; the
model is trained on the anomaly-free run plus the normal rows of the remaining
training portions; the last 20\% of every run is the test set (7{,}503 rows,
30.0\% anomalous).
\paragraph{SKAB, benchmark protocol.} The SKAB reference implementation trains
a model per run on its first 400 rows and tests on the rest. This protocol is
also reported, with the label-free threshold set from each run's own training
scores, so results can be read against the SKAB leaderboard.
\paragraph{SWaT, standard protocol.} Detectors are trained on the normal run
from row 21{,}601 to 85\% of its length; the last 15\% is held-out normal data
for the label-free threshold; the first third of the attack period is the
labelled validation set for the validation-MCC threshold. Threshold-free metrics
and label-free results are reported on the full attack period; both threshold
rules are also compared on its last two thirds, which neither rule has seen.

\subsection{Metrics and uncertainty}
AUROC and average precision (AUPRC) summarise ranking quality independent of
any threshold. At a threshold, precision, recall, F1, MCC, false-alarm rate on
normal rows, event recall (the fraction of contiguous anomalous segments with at
least one alarm) and median detection delay are reported. MCC is
\begin{equation}
\mathrm{MCC}=\frac{TP\cdot TN-FP\cdot FN}{\sqrt{(TP+FP)(TP+FN)(TN+FP)(TN+FN)}},
\end{equation}
which is zero for any predictor independent of the labels, including one that
flags every row; F1 for that predictor equals $2\pi/(1+\pi)$ for anomaly
prevalence $\pi$ (0.46 on the SKAB test set), which is why F1 alone can credit
a detector that separates nothing \citep{chicco2020}. Precision is bounded by
prevalence in the same way: with false-alarm rate $f$ and recall $r$,
$\mathrm{precision}=\pi r/(\pi r+(1-\pi)f)$, so the same detector looks much
worse in precision on SWaT ($\pi{=}0.121$) than on a test set built with
$\pi{=}0.315$, as in the original study.
Test rows are strongly autocorrelated, so 95\% confidence intervals are
obtained by resampling whole blocks with replacement: the 34 test files on SKAB
and non-overlapping one-hour blocks on SWaT. The same resamples are used for
every detector, so paired differences are valid. Isolation Forest is repeated
over ten seeds and the LSTM autoencoder over three seeds on SKAB.

\subsection{Drift and retraining}
The retraining study simulates deployment on SWaT. An initial model is fitted
on the first 50{,}000 normal rows after start-up (about 14 hours). The
remainder of the normal week and then the attack period are streamed in
one-hour batches; each batch is scored \emph{before} it can be used for
retraining. Four policies are compared: \emph{static} (never retrain);
\emph{drift-triggered} (retrain when the mean two-sample Kolmogorov--Smirnov
statistic between the batch and the current reference,
$\bar D=\frac1d\sum_j \sup_x|F^{\mathrm{ref}}_j(x)-F^{\mathrm{win}}_j(x)|$
\citep{massey1951}, exceeds the 95th percentile of its within-initial-window
values); \emph{periodic} (the same number of retrains as the triggered policy,
evenly spaced); and \emph{always} (after every batch). Each retrain uses the
last five batches, with rows selected in one of three ways: \emph{oracle}
(ground-truth normal rows, as in the original study), \emph{predicted-normal}
(rows the current model did not flag) or \emph{all} rows. After every fit the
threshold is the 99th percentile of the fit-window scores. Isolation Forest and
Elliptic Envelope are tested. The scaler is fitted once on the initial window.
The KS statistic is used instead of its $p$-value because, with thousands of
rows per batch, $p$-values vanish for negligible differences, whereas $\bar D$
is a bounded effect size.

\subsection{Resource measurement}
Each detector is fitted and timed in its own process inside a Docker container
limited to one CPU and 1\,GB of RAM on an ARM64 (Apple M4) host, with BLAS
threads limited to one. Reported are fit time, serialised size, load time,
single-reading latency over 1{,}000 calls (median, 95th and 99th percentiles),
batch throughput on 10{,}000 rows, and the child process's peak resident memory.
The cost of recomputing the rolling features for one new reading from a
30-row buffer is reported separately. This is a capped container, not an edge
device: a physical device's slower cores, thermal limits and I/O are not
represented, so latencies are lower bounds.

\section{Results}
\label{sec:results}

\subsection{SKAB under the pooled protocol}
Table~\ref{tab:skab} and Fig.~\ref{fig:auroc} (left) report SKAB under the pooled
protocol. Every detector is close to chance. The highest AUROC is Isolation
Forest's $0.575$ ($95\%$ CI $0.504$--$0.645$; mean $0.575\pm0.018$ over ten seeds),
and its CI barely excludes $0.5$; the CIs of the top five detectors overlap, and
the paired difference between the two best (Isolation Forest and COPOD) is
$0.012$ (CI $-0.066$ to $0.092$). The LSTM autoencoder is below chance
(AUROC $0.442$, three seeds within $\pm0.001$). With a calibrated threshold the
best MCC is $0.113$ (COPOD; CI $0.045$--$0.190$). The label-free rule gives MCC
within $\pm0.05$ of zero for every detector: thresholds learned from normal
data in this pooled setting do not separate faults.

\begin{table}[t]
\centering\small
\caption{SKAB, pooled protocol (test: last 20\% of each of 34 runs, 7{,}503 rows,
30.0\% anomalous). AUROC and MCC with 95\% block-bootstrap CIs over files. MCC
columns use the label-free (LF) or validation-MCC (VM) threshold; event recall
uses VM; 33 fault events.}
\label{tab:skab}
\begin{tabular}{lcccccc}
\toprule
Detector & AUROC [95\% CI] & AUPRC & MCC LF & MCC VM [95\% CI] & Event rec. VM \\
\midrule
Isolation Forest & 0.575 [0.504, 0.645] & 0.338 & 0.005 & 0.097 [$-$0.026, 0.233] & 0.94 \\
COPOD & 0.558 [0.507, 0.621] & 0.320 & $-$0.042 & 0.113 [0.045, 0.190] & 1.00 \\
ECOD & 0.535 [0.476, 0.601] & 0.315 & $-$0.016 & 0.066 [$-$0.010, 0.159] & 1.00 \\
Elliptic Envelope & 0.520 [0.459, 0.581] & 0.305 & $-$0.050 & 0.051 [$-$0.079, 0.182] & 0.45 \\
$z$-score & 0.513 [0.477, 0.554] & 0.307 & $-$0.051 & $-$0.032 [$-$0.102, 0.039] & 0.48 \\
HBOS & 0.510 [0.420, 0.602] & 0.315 & 0.041 & 0.012 [$-$0.116, 0.152] & 0.85 \\
One-Class SVM & 0.498 [0.463, 0.533] & 0.300 & $-$0.035 & $-$0.042 [$-$0.106, 0.009] & 0.73 \\
PCA-recon & 0.489 [0.455, 0.517] & 0.294 & $-$0.046 & $-$0.047 [$-$0.104, $-$0.005] & 0.76 \\
LOF & 0.489 [0.450, 0.533] & 0.299 & 0.007 & $-$0.021 [$-$0.073, 0.032] & 0.82 \\
LSTM-AE & 0.442 [0.392, 0.500] & 0.267 & $-$0.075 & $-$0.078 [$-$0.165, $-$0.007] & 0.27 \\
\bottomrule
\end{tabular}
\end{table}

\begin{figure}[t]
\centering
\includegraphics[width=\textwidth]{auroc_ci.png}
\caption{AUROC with 95\% block-bootstrap confidence intervals (blocks: test files
on SKAB, one-hour blocks on SWaT). Dashed line: chance. On SKAB (pooled) no
detector is clearly above chance; on SWaT (standard split) all ten are well above
chance and statistically indistinguishable from one another.}
\label{fig:auroc}
\end{figure}

\subsection{SWaT under the standard protocol}
Table~\ref{tab:swat} and Fig.~\ref{fig:auroc} (right) report SWaT. All ten
detectors rank attack rows well (AUROC $0.825$--$0.858$) and their CIs overlap
almost completely; the difference between the top two (ECOD and One-Class SVM) is
$0.005$ (CI $-0.033$ to $0.044$). Isolation Forest's AUROC over ten seeds is
$0.824\pm0.008$. The LSTM autoencoder ranks attacks as well as the classical
detectors (AUROC $0.852$, AUPRC $0.763$, the highest AUPRC).

The detectors separate once a threshold is applied. With the label-free rule,
which targets 1\% false alarms on held-out normal-week data, the $z$-score
baseline reaches MCC $0.716$ (CI $0.589$--$0.819$), One-Class SVM $0.660$ and
Isolation Forest $0.637$, with false-alarm rates of 1.4--3.8\% on the
attack period's normal rows and alarms in 63--86\% of the 35 attack segments.
LOF, PCA reconstruction and the LSTM autoencoder instead raise alarms on
68--73\% of the attack period's \emph{normal} rows, although their threshold
gave exactly 1\% on held-out normal-week data: their scores shift between the
normal week and the attack week independently of the attacks, so their
thresholds do not transfer. Elliptic Envelope shows the opposite failure, with
very few alarms (MCC $0.008$). When a labelled third of the attack period is
used to calibrate thresholds, MCC on the remaining two thirds rises for most
detectors (One-Class SVM $0.828$, $z$-score $0.827$, COPOD $0.802$), but not for
the LSTM autoencoder ($0.266$), LOF ($0.246$) or PCA reconstruction ($0.143$).

\begin{table}[t]
\centering\small
\caption{SWaT, standard protocol (test: full attack period, 449{,}919 rows,
12.1\% attacks, 35 attack segments). AUROC and MCC with 95\% block-bootstrap CIs
(one-hour blocks). LF: label-free threshold (99th percentile of held-out
normal-week scores); FAR: false-alarm rate on the attack period's normal rows;
VM: validation-MCC threshold calibrated on the first third, scored on the last
two thirds.}
\label{tab:swat}
\begin{tabular}{lcccccc}
\toprule
Detector & AUROC [95\% CI] & AUPRC & MCC LF [95\% CI] & FAR LF & Event rec. LF & MCC VM \\
\midrule
ECOD & 0.858 [0.779, 0.925] & 0.703 & 0.310 [0.196, 0.395] & 0.001 & 0.20 & 0.688 \\
One-Class SVM & 0.853 [0.772, 0.916] & 0.752 & 0.660 [0.506, 0.776] & 0.038 & 0.86 & 0.828 \\
COPOD & 0.852 [0.772, 0.920] & 0.728 & 0.584 [0.438, 0.683] & 0.001 & 0.17 & 0.802 \\
LSTM-AE & 0.852 [0.755, 0.917] & 0.763 & 0.134 [0.059, 0.209] & 0.726 & 0.89 & 0.266 \\
LOF & 0.852 [0.759, 0.916] & 0.744 & 0.183 [0.129, 0.250] & 0.683 & 1.00 & 0.246 \\
$z$-score & 0.847 [0.769, 0.907] & 0.738 & 0.716 [0.589, 0.819] & 0.014 & 0.86 & 0.827 \\
PCA-recon & 0.845 [0.756, 0.910] & 0.747 & 0.144 [0.073, 0.214] & 0.704 & 0.91 & 0.143 \\
Isolation Forest & 0.838 [0.747, 0.912] & 0.720 & 0.637 [0.480, 0.755] & 0.034 & 0.63 & 0.738 \\
Elliptic Envelope & 0.829 [0.737, 0.899] & 0.528 & 0.008 [$-$0.029, 0.056] & 0.014 & 0.74 & 0.643 \\
HBOS & 0.825 [0.734, 0.903] & 0.554 & 0.500 [0.333, 0.629] & 0.088 & 0.63 & 0.443 \\
\bottomrule
\end{tabular}
\end{table}

\subsection{How much protocol choices change the conclusions}
\label{sec:protocol}

\paragraph{Split definition (SWaT).} Scoring the same detectors with the same
code on the re-upload's layout (the split of the original study: train on the
normal rows between attacks, test on 31.5\% attacks) gives AUROC $0.840$--$0.934$
and reorders the detectors: HBOS ranks first there and last under the standard
split (Fig.~\ref{fig:protocol}, left). Under that layout the training data come
from the same five days as the test attacks, so the normal-week to attack-week
shift described above never appears in evaluation.

\paragraph{Threshold rule (SWaT).} The original study tuned one contamination
value on Isolation Forest and applied it to every model. On the same layout,
per-model validation-MCC thresholds (with the revised detector settings of
Section~\ref{sec:method}) raise the $z$-score's MCC from $0.180$ to
$0.844$, PCA reconstruction's from $0.002$ to $0.834$ and LOF's from $0.105$ to
$0.830$, while Elliptic Envelope, the original winner at $0.766$, moves only to
$0.810$. The original ranking therefore reflected mainly the shared threshold,
not the detectors.

\paragraph{Training protocol (SKAB).} Under the SKAB reference protocol (one
model per run, trained on its first 400 rows), the same detectors are far
better than under pooled training: mean per-run AUROC reaches $0.794$
(PCA reconstruction), $0.785$ (LOF) and $0.771$ (One-Class SVM), against
$0.489$--$0.575$ pooled (Fig.~\ref{fig:protocol}, right; Table~\ref{tab:skabbench}).
The ranking inverts: Isolation Forest, the best pooled detector, is among the
weakest per run, and ECOD and COPOD fall to chance because, with only 400 training
rows and a 99th-percentile threshold, they almost never alarm. Of the first 400
rows, 2.1\% are labelled anomalous, so this protocol trains on slightly
contaminated data, as in the reference implementation. Absolute numbers are not
directly comparable with the SKAB leaderboard (best F1 $0.78$ at a 13.6\% false-alarm
rate for a convolutional autoencoder), because features, thresholds and post-processing
differ; detectors here reach similar F1 ($0.72$--$0.73$) only at much higher
false-alarm rates (47--71\%), while Elliptic Envelope offers a more balanced
operating point (F1 $0.644$, false alarms 20.2\%, MCC $0.361$).

\paragraph{Threshold rule (SKAB).} On the pooled SKAB test set, Isolation
Forest's MCC is $0.005$ with the label-free rule, $0.097$ with validation-MCC and
$0.118$ with validation-F1; the F1 rule buys its MCC with recall $0.71$ at
higher false-alarm rates. Differences of this size are within the CIs.

\paragraph{Feature choices.} Removing the elapsed-time feature and computing
features before dropping fault rows changed SKAB AUROC by at most $0.02$
(Isolation Forest $0.591$ original pipeline, $0.555$ with elapsed time, $0.575$
revised), well within the CIs; these corrections matter for validity, not for
the numbers.

\begin{figure}[t]
\centering
\includegraphics[width=\textwidth]{protocol_sensitivity.png}
\caption{Same detectors, same code, different protocol. Left: SWaT AUROC under the
re-upload layout used in the original study versus the standard split. Right: SKAB
AUROC with one model trained on all runs (pooled) versus one model per run trained
on its first 400 rows (SKAB reference protocol; mean over runs).}
\label{fig:protocol}
\end{figure}

\begin{table}[t]
\centering\small
\caption{SKAB reference protocol (per run: first 400 rows train, remainder test;
label-free threshold from each run's training scores). FAR and MAR as defined by the
SKAB benchmark (false-alarm and missed-alarm rates).}
\label{tab:skabbench}
\begin{tabular}{lccccc}
\toprule
Detector & F1 & FAR & MAR & MCC & Mean run AUROC \\
\midrule
Elliptic Envelope & 0.644 & 0.202 & 0.445 & 0.361 & 0.749 \\
$z$-score & 0.722 & 0.468 & 0.213 & 0.331 & 0.766 \\
LOF & 0.723 & 0.556 & 0.168 & 0.302 & 0.785 \\
One-Class SVM & 0.732 & 0.712 & 0.077 & 0.277 & 0.771 \\
PCA-recon & 0.731 & 0.707 & 0.081 & 0.276 & 0.794 \\
HBOS & 0.603 & 0.306 & 0.457 & 0.239 & 0.674 \\
Isolation Forest & 0.470 & 0.201 & 0.641 & 0.174 & 0.664 \\
COPOD & 0.044 & 0.024 & 0.977 & $-$0.004 & 0.548 \\
ECOD & 0.027 & 0.022 & 0.986 & $-$0.029 & 0.493 \\
\bottomrule
\end{tabular}
\end{table}

\section{Drift and Retraining}
\label{sec:drift}

\paragraph{Does normal operation drift?} Streaming SWaT's genuine normal week
through a static Isolation Forest fitted on its first 14 hours, the share of
normal readings that raise an alarm, after a transient spike in the first day,
stays near the design rate of 1\% for about two days and then jumps, at roughly
hour 48, to $40$--$94\%$ for the rest of the week (mean $12\%$ over the first 24
streamed hours, $76\%$ over the last 24; Fig.~\ref{fig:drift}, left). The change
is a step, consistent with a change in operating regime, rather than a slow trend. The
Elliptic Envelope fitted on the same window does not show this: its false-alarm
rate over the normal week is $0.9\%$. Drift is therefore real, but whether it
matters depends on the detector. Isolation Forest operates on all 120 raw-scale
features, whereas the Elliptic Envelope operates on ten whitened principal
components; a plausible, untested explanation is that slow level changes in
individual sensors fall mostly outside those components.

\paragraph{Retraining policies.} Table~\ref{tab:drift} and Fig.~\ref{fig:drift}
(right) compare policies over the attack period after the normal week. With
retraining data selected by ground-truth labels (as in the original study),
retraining helps Isolation Forest substantially: MCC rises from $0.167$
(static, which flags 73\% of normal rows) to $0.488$ when triggered by the KS
statistic (74 retrains), $0.456$ when the same number of retrains is evenly
spaced, and $0.537$ when retraining after every batch. Drift triggering is thus
about as good as a schedule with the same budget; it is not clearly better. For
the Elliptic Envelope, no policy improves on the static model (MCC $0.723$ static,
$0.709$--$0.734$ with retraining).

\paragraph{Without labels, retraining can fail.} A deployed detector does not
know which recent rows are normal. When Isolation Forest retrains on the rows it
did \emph{not} flag, it enters a self-reinforcing failure: once it flags most rows,
too few remain to retrain on, and it ends up flagging everything (MCC $0.000$ to
$0.025$, false-alarm rate 99.5--100\%). Retraining on all recent rows avoids that
collapse but absorbs the attacks into the model of normal behaviour, roughly
halving recall ($0.30$--$0.51$) and giving MCC $0.188$--$0.303$. The Elliptic
Envelope tolerates predicted-normal selection (MCC $0.683$--$0.734$) but loses
more than half its skill when retraining on all rows ($0.332$--$0.418$). The
benefit of adaptation reported in the original study therefore depended on
labels that a deployment would not have.

\paragraph{Cost.} On the host, an Isolation Forest retrain on five hours of data
took about $0.08$--$0.1$\,s and an Elliptic Envelope retrain about $1.1$\,s. The KS
drift check (120 features against a reference of up to 50{,}000 rows) took
$0.22$--$0.46$\,s per hourly batch in the label-selected and all-rows runs, and up
to $2$\,s when the reference grew large. For Isolation Forest, checking for drift
therefore costs more than simply retraining every hour, and every-hour
retraining gave the best result. A drift trigger pays off only when retraining
is expensive relative to the check, as for the Elliptic Envelope.

\begin{table}[t]
\centering\small
\caption{Retraining policies on SWaT (stream: rest of the normal week, 119 hourly
batches, then the attack period, 125 batches). Normal FAR: false-alarm rate over the
streamed normal week. Attack-period metrics are point-wise; event recall over 35
segments. Retrain CPU: total seconds spent refitting (host).}
\label{tab:drift}
\begin{tabular}{llrcccccr}
\toprule
Detector & Policy & Retrains & Normal FAR & MCC & Recall & FAR & Event rec. & CPU s \\
\midrule
\multicolumn{9}{l}{\emph{Isolation Forest}}\\
 & static & 0 & 0.431 & 0.167 & 0.949 & 0.730 & 1.00 & 0.0 \\
 & KS, oracle & 74 & 0.119 & 0.488 & 0.695 & 0.115 & 0.83 & 7.0 \\
 & periodic, oracle & 80 & 0.187 & 0.456 & 0.716 & 0.144 & 0.83 & 6.4 \\
 & always, oracle & 239 & 0.105 & 0.537 & 0.687 & 0.086 & 0.80 & 19.6 \\
 & KS, pred.-normal & 52 & 0.565 & 0.025 & 1.000 & 0.995 & 1.00 & 4.1 \\
 & periodic, pred.-normal & 10 & 0.877 & 0.000 & 1.000 & 1.000 & 1.00 & 0.5 \\
 & always, pred.-normal & 52 & 0.721 & 0.001 & 1.000 & 1.000 & 1.00 & 3.3 \\
 & KS, all rows & 79 & 0.119 & 0.188 & 0.305 & 0.110 & 0.86 & 7.2 \\
 & periodic, all rows & 81 & 0.187 & 0.303 & 0.505 & 0.143 & 0.86 & 7.8 \\
 & always, all rows & 244 & 0.105 & 0.226 & 0.298 & 0.083 & 0.83 & 22.7 \\
\multicolumn{9}{l}{\emph{Elliptic Envelope}}\\
 & static & 0 & 0.009 & 0.723 & 0.634 & 0.011 & 0.80 & 0.0 \\
 & KS, oracle & 74 & 0.011 & 0.709 & 0.598 & 0.008 & 0.74 & 79.6 \\
 & periodic, oracle & 80 & 0.021 & 0.717 & 0.615 & 0.010 & 0.74 & 95.3 \\
 & always, oracle & 239 & 0.011 & 0.711 & 0.607 & 0.010 & 0.74 & 283.0 \\
 & KS, pred.-normal & 73 & 0.014 & 0.725 & 0.643 & 0.012 & 0.69 & 85.9 \\
 & periodic, pred.-normal & 80 & 0.025 & 0.734 & 0.667 & 0.014 & 0.69 & 84.1 \\
 & always, pred.-normal & 238 & 0.016 & 0.683 & 0.653 & 0.024 & 0.77 & 234.7 \\
 & KS, all rows & 79 & 0.011 & 0.332 & 0.180 & 0.008 & 0.77 & 82.4 \\
 & periodic, all rows & 81 & 0.021 & 0.418 & 0.260 & 0.010 & 0.80 & 90.0 \\
 & always, all rows & 244 & 0.011 & 0.335 & 0.189 & 0.010 & 0.80 & 279.1 \\
\bottomrule
\end{tabular}
\end{table}

\begin{figure}[t]
\centering
\includegraphics[width=\textwidth]{drift_policies.png}
\caption{Left: false-alarm rate of a static Isolation Forest over SWaT's genuine
normal week (threshold set to flag 1\% of its training window). Middle and right: detection
skill (MCC) over the attack period for each retraining policy, coloured by how
retraining data are selected; grey is the static model. Selecting with ground-truth
labels (blue) is what makes retraining help Isolation Forest; label-free selection
(orange, green) does not. The static Elliptic Envelope is already as good as any
adaptive variant.}
\label{fig:drift}
\end{figure}

\paragraph{A note on the Elliptic Envelope's instability.} In this streaming
setting the static Elliptic Envelope, fitted on 14 hours, reaches MCC $0.723$,
yet in Table~\ref{tab:swat}, fitted on the full normal week with its threshold
set on the week's last 15\%, its label-free MCC is $0.008$. The two runs differ
only in the training window and threshold-calibration data. One consistent reading
is that drift within the normal week inflates the end-of-week scores used for
calibration, producing a threshold too high to fire. Either way, a detector
whose performance swings this much with the calibration window is fragile in
deployment, a property a single benchmark number hides.

%%EDGE%%

\section{Discussion}
\label{sec:discussion}

\paragraph{Ranking quality is not the bottleneck; thresholds and protocols are.}
On SWaT's standard split, ten detectors that span statistical, distance,
density, tree, copula and recurrent families are statistically indistinguishable
in AUROC. This agrees with large benchmark studies that find no consistent winner
and strong simple baselines \citep{audibert2022,sarfraz2024}. What separates
them in use is whether a threshold set on normal data still works later: the
$z$-score, One-Class SVM and Isolation Forest kept false alarms near their design
rate on the attack week's normal data, while LOF, PCA reconstruction and the LSTM
autoencoder did not. A detector comparison that reports only threshold-free
ranking, or only thresholds tuned on test-period labels, hides exactly the property
an operator cares about.

\paragraph{Why the original conclusions changed.} Three protocol choices in the
original study each favoured a particular answer. Training on the attack period's
own normal rows removed the normal-week to attack-week shift from the evaluation.
A contamination value tuned on Isolation Forest and shared by all detectors made
detectors with differently scaled scores look poor. Retraining on
ground-truth-normal rows made adaptation look free of risk. None of these
choices is unusual, and each was documented in the code. Their effect became
visible only when the same detectors were rerun under alternative protocols.
Studies that compare detectors should therefore report at least one alternative
threshold rule and state where their training and test data come from in time.

\paragraph{What adaptation needs.} Normal operation does drift, and some
detectors are much more sensitive to it than others. Where retraining helps, it
helps only if the retraining data are clean; with label-free selection Isolation
Forest either collapsed into flagging everything or learned the attacks as
normal. In practice this argues for one of three things: retraining data vetted
by an operator; a slow-changing detector (here, the Elliptic Envelope on whitened
components), whose static form was as good as any adaptive variant; or adaptation
of the threshold rather than the model. For cheap models such as Isolation Forest, a
drift test can cost more than the retrain it is meant to save.

\paragraph{Cost.} Single-reading latencies span three orders of magnitude, from
microseconds ($z$-score) to over 100\,ms (ECOD and COPOD in PyOD, which store and
rescore against their training data), and model sizes from 1\,KB to 39\,MB. The
batch-amortised figures used in the original study (around $0.002$\,ms per row for
Isolation Forest) understate the real single-reading cost (about $2.3$--$2.8$\,ms)
by three orders of magnitude, because per-call overhead dominates. For a monitor that
scores one reading per second, all detectors except possibly ECOD and COPOD fit
comfortably; the deciding constraints are memory and threshold stability, not
speed.

\section{Threats to Validity and Limitations}
\label{sec:limits}

\begin{itemize}
\item \textbf{Dataset copy.} SWaT results come from a public re-upload
reconstructed to the standard split; seven channels are blank in its normal run, so 24
of the 25 continuous sensors were used and actuators were excluded. The
reconstruction reproduces the published row counts, but the results must be
confirmed on the official iTrust release before they are cited as SWaT results.
\item \textbf{Two water systems.} Both datasets are water processes with
physically induced faults or scripted attacks; rotating machinery, electrical
faults and progressive wear are not covered. HAI \citep{shin2020} and WADI
\citep{ahmed2017} are natural next datasets.
\item \textbf{Detector configurations.} Hyperparameters other than thresholds were
fixed at common defaults; One-Class SVM and LOF were trained on a 20{,}000-row
subsample on SWaT; the LSTM autoencoder is small and was run with one seed on
SWaT. Larger or tuned deep models were not tested.
\item \textbf{Metrics.} Event recall and median delay are reported, but
range-aware metrics such as VUS-PR \citep{paparrizos2022} were not computed.
Bootstrap blocks (files, hours) reduce but do not remove dependence between blocks.
\item \textbf{Drift design.} One stream, one window length (five hours), one KS
threshold rule and two detectors were studied; ADWIN-style detectors and threshold-only
adaptation were not compared. Policies are compared without confidence intervals
because the stream is a single realisation.
\item \textbf{Resource measurements.} A Docker container capped at one CPU and
1\,GB on an Apple M4 host is not a physical edge device: slower cores, thermal
limits and real I/O are absent, so latencies are lower bounds. Peak memory is per
process and includes loading the dataset. Energy was not measured.
\item \textbf{SKAB protocols.} The SKAB reference protocol trains on rows of which
2.1\% are labelled anomalous, and its label-free threshold is set from each
run's own training scores; results are not directly comparable with the SKAB
leaderboard, whose features, thresholds and post-processing differ.
\item \textbf{Search scope.} The literature search used a general web search and
publisher pages rather than database-level systematic queries; statements that a
comparison was not found are limited to that search.
\end{itemize}

\section{Conclusion}
\label{sec:conclusion}

Re-examining a six-detector comparison on SKAB and SWaT under documented,
standard protocols answered the four research questions as follows.
(RQ1) On SWaT's standard split, ten lightweight and deep detectors rank attacks
equally well (AUROC $0.83$--$0.86$, overlapping intervals); on SKAB with pooled
training, none is reliably better than chance. (RQ2) Protocol choices change the
conclusions more than detector choice does: the SWaT split definition reorders the
detectors, a shared contamination value had made the $z$-score look useless (MCC
$0.18$ against $0.84$ with its own threshold), and SKAB's per-run protocol inverts the
pooled ranking. Under thresholds set only from normal data, the simple $z$-score,
One-Class SVM and Isolation Forest were the most dependable on SWaT (MCC
$0.64$--$0.72$). (RQ3) Normal operation drifts enough to break a static Isolation
Forest; retraining restores it only when the retraining data are known to be
normal, drift-triggered retraining was no better than a schedule with the same
budget, and label-free retraining failed. (RQ4) Single-reading costs range from
microseconds to over 100\,ms and model sizes from kilobytes to tens of megabytes
under a one-CPU, 1\,GB cap.

The original version of this study reported dataset-specific winners, a large
benefit from drift adaptation, and sub-millisecond inference; none of these
survived re-examination. What did survive is a clearer and more useful message for
practitioners: choose the detector for its threshold stability and retraining
needs rather than its benchmark rank, and report the protocol in enough detail
that a reader can tell which of these the numbers measure. Next steps are
confirmation on the official SWaT release, a third dataset from a different
machine class, threshold-only adaptation, and measurement on physical edge
hardware.

\section*{Data and Code Availability}
Code, configuration and all result files are available in the project repository
(\texttt{experiments/revision/}, \texttt{results/revised/}). SKAB is public
\citep{skab}. SWaT must be requested from iTrust, SUTD; the re-upload used and the
reconstruction of the standard split are documented in the code. Software used:
scikit-learn \citep{pedregosa2011}, PyOD \citep{zhao2019} and PyTorch.

\section*{Disclosure of Interests}
The authors have no competing interests to declare.

\begin{thebibliography}{99}
\small

\bibitem[Ahmad et~al.(2017)]{ahmad2017} Ahmad, S., Lavin, A., Purdy, S., \& Agha, Z. (2017). Unsupervised real-time anomaly detection for streaming data. \emph{Neurocomputing}, 262, 134--147. doi:10.1016/j.neucom.2017.04.070

\bibitem[Ahmed et~al.(2017)]{ahmed2017} Ahmed, C.~M., Palleti, V.~R., \& Mathur, A.~P. (2017). WADI: A water distribution testbed for research in the design of secure cyber physical systems. In \emph{Proc.\ 3rd Int.\ Workshop on Cyber-Physical Systems for Smart Water Networks (CySWater)} (pp.\ 25--28). doi:10.1145/3055366.3055375

\bibitem[Audibert et~al.(2020)]{audibert2020} Audibert, J., Michiardi, P., Guyard, F., Marti, S., \& Zuluaga, M.~A. (2020). USAD: Unsupervised anomaly detection on multivariate time series. In \emph{Proc.\ 26th ACM SIGKDD Int.\ Conf.\ on Knowledge Discovery \& Data Mining}. doi:10.1145/3394486.3403392

\bibitem[Audibert et~al.(2022)]{audibert2022} Audibert, J., Michiardi, P., Guyard, F., Marti, S., \& Zuluaga, M.~A. (2022). Do deep neural networks contribute to multivariate time series anomaly detection? \emph{Pattern Recognition}, 132, 108945. doi:10.1016/j.patcog.2022.108945

\bibitem[Banbury et~al.(2021)]{banbury2021} Banbury, C., et~al. (2021). MLPerf Tiny benchmark. In \emph{Proc.\ NeurIPS Datasets and Benchmarks Track}. arXiv:2106.07597

\bibitem[Bifet \& Gavald\`a(2007)]{bifet2007} Bifet, A., \& Gavald\`a, R. (2007). Learning from time-changing data with adaptive windowing. In \emph{Proc.\ SIAM Int.\ Conf.\ on Data Mining} (pp.\ 443--448). doi:10.1137/1.9781611972771.42

\bibitem[Breunig et~al.(2000)]{breunig2000} Breunig, M.~M., Kriegel, H.-P., Ng, R.~T., \& Sander, J. (2000). LOF: Identifying density-based local outliers. In \emph{Proc.\ ACM SIGMOD Int.\ Conf.\ on Management of Data} (pp.\ 93--104). doi:10.1145/342009.335388

\bibitem[Chandola et~al.(2009)]{chandola2009} Chandola, V., Banerjee, A., \& Kumar, V. (2009). Anomaly detection: A survey. \emph{ACM Computing Surveys}, 41(3), 1--58. doi:10.1145/1541880.1541882

\bibitem[Chicco \& Jurman(2020)]{chicco2020} Chicco, D., \& Jurman, G. (2020). The advantages of the Matthews correlation coefficient (MCC) over F1 score and accuracy in binary classification evaluation. \emph{BMC Genomics}, 21, 6. doi:10.1186/s12864-019-6413-7

\bibitem[Deng \& Hooi(2021)]{deng2021} Deng, A., \& Hooi, B. (2021). Graph neural network-based anomaly detection in multivariate time series. In \emph{Proc.\ AAAI Conf.\ on Artificial Intelligence}, 35(5), 4027--4035. doi:10.1609/aaai.v35i5.16523

\bibitem[Gama et~al.(2014)]{gama2014} Gama, J., \v{Z}liobait\.{e}, I., Bifet, A., Pechenizkiy, M., \& Bouchachia, A. (2014). A survey on concept drift adaptation. \emph{ACM Computing Surveys}, 46(4), 1--37. doi:10.1145/2523813

\bibitem[Garg et~al.(2021)]{garg2021} Garg, A., Zhang, W., Samaran, J., Savitha, R., \& Foo, C.-S. (2021). An evaluation of anomaly detection and diagnosis in multivariate time series. \emph{IEEE Transactions on Neural Networks and Learning Systems}. doi:10.1109/TNNLS.2021.3105827

\bibitem[Goh et~al.(2016)]{goh2016} Goh, J., Adepu, S., Junejo, K.~N., \& Mathur, A. (2016). A dataset to support research in the design of secure water treatment systems. In \emph{Proc.\ 11th Int.\ Conf.\ on Critical Information Infrastructures Security (CRITIS)} (pp.\ 88--99). doi:10.1007/978-3-319-71368-7\_8

\bibitem[Goldstein \& Uchida(2016)]{goldstein2016} Goldstein, M., \& Uchida, S. (2016). A comparative evaluation of unsupervised anomaly detection algorithms for multivariate data. \emph{PLoS ONE}, 11(4), e0152173. doi:10.1371/journal.pone.0152173

\bibitem[Han et~al.(2022)]{han2022} Han, S., Hu, X., Huang, H., Jiang, M., \& Zhao, Y. (2022). ADBench: Anomaly detection benchmark. In \emph{Advances in Neural Information Processing Systems 35 (Datasets and Benchmarks Track)}. arXiv:2206.09426

\bibitem[Hundman et~al.(2018)]{hundman2018} Hundman, K., Constantinou, V., Laporte, C., Colwell, I., \& Soderstrom, T. (2018). Detecting spacecraft anomalies using LSTMs and nonparametric dynamic thresholding. In \emph{Proc.\ 24th ACM SIGKDD Int.\ Conf.\ on Knowledge Discovery \& Data Mining}. doi:10.1145/3219819.3219845

\bibitem[Katser \& Kozitsin(2020)]{skab} Katser, I.~D., \& Kozitsin, V.~O. (2020). \emph{Skoltech Anomaly Benchmark (SKAB)} [Data set]. Kaggle. doi:10.34740/KAGGLE/DSV/1693952

\bibitem[Kim et~al.(2022)]{kim2022} Kim, S., Choi, K., Choi, H.-S., Lee, B., \& Yoon, S. (2022). Towards a rigorous evaluation of time-series anomaly detection. In \emph{Proc.\ AAAI Conf.\ on Artificial Intelligence}, 36(7), 7194--7201. doi:10.1609/aaai.v36i7.20680

\bibitem[Li et~al.(2019)]{li2019madgan} Li, D., Chen, D., Jin, B., Shi, L., Goh, J., \& Ng, S.-K. (2019). MAD-GAN: Multivariate anomaly detection for time series data with generative adversarial networks. In \emph{Proc.\ Int.\ Conf.\ on Artificial Neural Networks (ICANN)}. arXiv:1809.04758

\bibitem[Li et~al.(2023)]{li2022ecod} Li, Z., Zhao, Y., Hu, X., Botta, N., Ionescu, C., \& Chen, G.~H. (2023). ECOD: Unsupervised outlier detection using empirical cumulative distribution functions. \emph{IEEE Transactions on Knowledge and Data Engineering}, 35(12), 12181--12193. doi:10.1109/TKDE.2022.3159580

\bibitem[Liu \& Paparrizos(2024)]{liu2024} Liu, Q., \& Paparrizos, J. (2024). The elephant in the room: Towards a reliable time-series anomaly detection benchmark. In \emph{Advances in Neural Information Processing Systems 37 (Datasets and Benchmarks Track)} (pp.\ 108231--108261).

\bibitem[Liu et~al.(2008)]{liu2008} Liu, F.~T., Ting, K.~M., \& Zhou, Z.-H. (2008). Isolation forest. In \emph{Proc.\ 8th IEEE Int.\ Conf.\ on Data Mining} (pp.\ 413--422). doi:10.1109/ICDM.2008.17

\bibitem[Lu et~al.(2019)]{lu2018} Lu, J., Liu, A., Dong, F., Gu, F., Gama, J., \& Zhang, G. (2019). Learning under concept drift: A review. \emph{IEEE Transactions on Knowledge and Data Engineering}, 31(12), 2346--2363. doi:10.1109/TKDE.2018.2876857

\bibitem[Malhotra et~al.(2016)]{malhotra2016} Malhotra, P., Ramakrishnan, A., Anand, G., Vig, L., Agarwal, P., \& Shroff, G. (2016). LSTM-based encoder-decoder for multi-sensor anomaly detection. \emph{ICML Anomaly Detection Workshop}. arXiv:1607.00148

\bibitem[Massey(1951)]{massey1951} Massey, F.~J. (1951). The Kolmogorov-Smirnov test for goodness of fit. \emph{Journal of the American Statistical Association}, 46(253), 68--78. doi:10.1080/01621459.1951.10500769

\bibitem[Mathur \& Tippenhauer(2016)]{mathur2016} Mathur, A.~P., \& Tippenhauer, N.~O. (2016). SWaT: A water treatment testbed for research and training on ICS security. In \emph{Proc.\ Int.\ Workshop on Cyber-physical Systems for Smart Water Networks (CySWater)} (pp.\ 31--36). doi:10.1109/CySWater.2016.7469060

\bibitem[Paparrizos et~al.(2022)]{paparrizos2022} Paparrizos, J., Boniol, P., Palpanas, T., Tsay, R.~S., Elmore, A., \& Franklin, M.~J. (2022). Volume under the surface: A new accuracy evaluation measure for time-series anomaly detection. \emph{Proceedings of the VLDB Endowment}, 15(11), 2774--2787. doi:10.14778/3551793.3551830

\bibitem[Pedregosa et~al.(2011)]{pedregosa2011} Pedregosa, F., et~al. (2011). Scikit-learn: Machine learning in Python. \emph{Journal of Machine Learning Research}, 12, 2825--2830.

\bibitem[Rabanser et~al.(2019)]{rabanser2019} Rabanser, S., G\"unnemann, S., \& Lipton, Z. (2019). Failing loudly: An empirical study of methods for detecting dataset shift. In \emph{Advances in Neural Information Processing Systems 32}. arXiv:1810.11953

\bibitem[Rousseeuw \& Van~Driessen(1999)]{rousseeuw1999} Rousseeuw, P.~J., \& Van~Driessen, K. (1999). A fast algorithm for the minimum covariance determinant estimator. \emph{Technometrics}, 41(3), 212--223. doi:10.1080/00401706.1999.10485670

\bibitem[Sarfraz et~al.(2024)]{sarfraz2024} Sarfraz, M.~S., Chen, M.-Y., Layer, L., Peng, K., \& Koulakis, M. (2024). Position: Quo vadis, unsupervised time series anomaly detection? In \emph{Proc.\ 41st Int.\ Conf.\ on Machine Learning}, PMLR 235, 43461--43476.

\bibitem[Schmidl et~al.(2022)]{schmidl2022} Schmidl, S., Wenig, P., \& Papenbrock, T. (2022). Anomaly detection in time series: A comprehensive evaluation. \emph{Proceedings of the VLDB Endowment}, 15(9), 1779--1797. doi:10.14778/3538598.3538602

\bibitem[Sch\"olkopf et~al.(2001)]{scholkopf2001} Sch\"olkopf, B., Platt, J.~C., Shawe-Taylor, J., Smola, A.~J., \& Williamson, R.~C. (2001). Estimating the support of a high-dimensional distribution. \emph{Neural Computation}, 13(7), 1443--1471. doi:10.1162/089976601750264965

\bibitem[Shin et~al.(2020)]{shin2020} Shin, H.-K., Lee, W., Yun, J.-H., \& Kim, H. (2020). HAI 1.0: HIL-based augmented ICS security dataset. In \emph{Proc.\ 13th USENIX Workshop on Cyber Security Experimentation and Test (CSET)}.

\bibitem[Shyu et~al.(2003)]{shyu2003} Shyu, M.-L., Chen, S.-C., Sarinnapakorn, K., \& Chang, L. (2003). A novel anomaly detection scheme based on principal component classifier. In \emph{Proc.\ IEEE ICDM Workshop on Foundations and New Directions of Data Mining} (pp.\ 172--179).

\bibitem[Tatbul et~al.(2018)]{tatbul2018} Tatbul, N., Lee, T.~J., Zdonik, S., Alam, M., \& Gottschlich, J. (2018). Precision and recall for time series. In \emph{Advances in Neural Information Processing Systems 31}. arXiv:1803.03639

\bibitem[Tuli et~al.(2022)]{tuli2022} Tuli, S., Casale, G., \& Jennings, N.~R. (2022). TranAD: Deep transformer networks for anomaly detection in multivariate time series data. \emph{Proceedings of the VLDB Endowment}, 15(6), 1201--1214. doi:10.14778/3514061.3514067

\bibitem[Wu \& Keogh(2023)]{wu2023} Wu, R., \& Keogh, E. (2023). Current time series anomaly detection benchmarks are flawed and are creating the illusion of progress. \emph{IEEE Transactions on Knowledge and Data Engineering}, 35(3), 2421--2429. doi:10.1109/TKDE.2021.3112126

\bibitem[Zhao et~al.(2019)]{zhao2019} Zhao, Y., Nasrullah, Z., \& Li, Z. (2019). PyOD: A Python toolbox for scalable outlier detection. \emph{Journal of Machine Learning Research}, 20(96), 1--7.

\end{thebibliography}

\end{document}
